\documentclass[10pt,a4paper]{article}
\usepackage[utf8]{inputenc}
\usepackage[a4paper,left=3cm,right=2cm]{geometry}
\usepackage{amsmath}
\usepackage[thmmarks, amsmath, thref]{ntheorem}
\usepackage{amsfonts}
\usepackage{amssymb}
\usepackage{fancyhdr}
\pagestyle{fancy}
\usepackage{anysize}
\usepackage{graphicx}
\usepackage{float}
\usepackage{hyperref}
\usepackage{enumitem}
\usepackage{verbatim}
\newtheorem{theorem}{Theorem}
\theoremsymbol{\ensuremath{\blacksquare}}
\newtheorem*{solution}{Solución:}

\usepackage{xcolor}

% Margenes y formalidades

\textwidth 6.25in % Margen Derecho y ancho del texto
\textheight 8.5in
\oddsidemargin 0in % Margen Izquierdo
\headheight 0in % Largo del texto
\leftmargin 1in
\parindent 0em % Salto entre la indentación (Sangría)https://www.overleaf.com/4519127669sdpptkhvsbfn
\topmargin -0.5in % Margen Superior
\parskip 0.5ex % Salto entre parrafos
\baselineskip 1pt



\lhead{
\includegraphics[scale=0.08]{Figuras/logoUSM.png}} % texto izquierda de la cabecera
%\chead{}                                      % texto centro de la cabecera
\rhead{Universidad Técnica Federico Santa María\\Ingeniería Comercial } % número de página a la derecha
\renewcommand{\headrulewidth}{0.4pt} % grosor de la línea de la cabecera
\renewcommand{\footrulewidth}{0.4pt} % grosor de la línea del pie

\begin{document}
\vspace{5cm}



\begin{center}
\begin{large}
\textbf{\\RECUPERATIVO CERTAMEN 1}
\end{large}\\
\end{center}

\begin{center}
\begin{large}
\textbf{Econometría ICS-294}
\end{large}\\
\vspace{0.2cm}

\\
%Ayudante: 
\end{center}

\textbf{Instrucciones:} 

\begin{itemize}

    \item Usted dispone de 90 minutos  para intentar los 100 puntos del certamen.
    \item Cuide la presentación y redacción de sus preguntas.
    \item No puede utilizar celular, computador o tablet, pero sí una calculadora TI.
     
\end{itemize}

\begin{enumerate}

\item  (10 puntos) Josefina se encuentra trabajando para una importante constructora nacional en el \'area de finanzas, el jefe de
proveedores le comenta que la media de cuentas por cobrar es de 220.000.000 de pesos. Josefina tiene una  muestra aleatoria de 500 facturas con media igual a 195.000.000 de pesos y una desviaci\'on
est\'andar de 1.000.000. Con un 90\% de confianza ¿es correcta la afirmación del jefe?


%     $$\color{blue}{\bar{x} \pm z_{\frac{\alpha}{2}}\frac{\sigma}{\sqrt{n}}}$$

% {\color{blue}$[194.926.209,8 ; 195.073.790,2]$ (6 puntos)}


%    {\color{blue}  El valor que comenta el jefe es 220.000.000 el cual no está en el intervalo de confianza calculado, así que  90\% de confianza la afirmación es incorrecta. (4 puntos)}



    \item (25 puntos) Considere el siguiente modelo de regresión lineal con los supuestos estándares:

    \[
    Y_i = \beta_0 + \beta_1 X_{1i} + \beta_2 X_{2i} + u_i
    \]
    
    Se dispone de la siguiente información:
    
    \[
    X'X = \begin{pmatrix}
    33 & 0 & 0 \\
    0 & 40 & 20 \\
    0 & 20 & 60
    \end{pmatrix}, \quad
    X'Y = \begin{pmatrix}
    132 \\
    24 \\
    92
    \end{pmatrix}
    \]
    
    \[
    \sum_{i=1}^{N} Y_i^2 = 678, \quad \sum_{i=1}^{N} (Y_i - \bar{Y})^2 = 150
    \]

%\begin{comment}
    \begin{enumerate}
        \item (5 puntos) ¿Cuál es el tamaño de la muestra?\\
    
        \item (12 puntos) Calcule los estimadores de MCO del modelo anterior.\\
        \textbf{Hint:} \( Z = \begin{pmatrix} A &{\bf 0} \\ {\bf 0} & B \end{pmatrix} \) con \( A \) y \( B \) matrices cuadradas, entonces:\[Z^{-1} = \begin{pmatrix} A^{-1} & {\bf 0} \\{\bf 0} & B^{-1} \end{pmatrix}\]

        \item (4 puntos) Indique la ecuación del modelo predicho e interprete.

        \item (4 puntos) ¿Qué debería suceder con $R^2$ si agregamos más variables independientes a la medición?
        
        % \item ¿Qué valor alcanza el $R^2$?\footnote{Hint: Recuerde que $SSE = \sum_{i=1}^N (\hat{Y_i} - \bar{Y})^2$ que de manera matricial queda como $SSE=\hat{Y}^T Y - N \bar{Y^2}$}
        %regresion lineal multiple
    \end{enumerate}
%\end{comment}

\begin{comment}
    \begin{enumerate}
        \item ¿Cuál es el tamaño de la muestra?\\

        \textbf{Respuesta:} El elemento (1, 1) de la matriz $X'X$ (sin expresarla en desviaciones) denota precisamente el tamaño muestral, por lo que éste es de 33 observaciones.\\
    
        \item Calcule los estimadores MELI del modelo anterior.\\
        \textbf{Hint:} \( Z = \begin{pmatrix} A & 0 \\ 0 & B \end{pmatrix} \) con \( A \) y \( B \) matrices cuadradas, entonces:\[Z^{-1} = \begin{pmatrix} A^{-1} & 0 \\ 0 & B^{-1} \end{pmatrix}\]

        \textbf{Respuesta:}
        Sabemos que los estimadores de mínimos cuadrados ordinarios (MCO) están dados por:
    
        \[
        \hat{\beta} = (X'X)^{-1} X'Y
        \]
    
        Calculando los estimadores:
    
        \[
        \hat{\beta} = \begin{pmatrix}
        33 & 0 & 0 \\
        0 & 40 & 20 \\
        0 & 20 & 60
        \end{pmatrix}^{-1}
        \begin{pmatrix}
        132 \\
        24 \\
        92
        \end{pmatrix}
        \]
        
        \[
        \hat{\beta} = \begin{pmatrix}
        \frac{1}{33} & 0 & 0 \\
        0 & 0.03 & -0.01 \\
        0 & -0.01 & 0.02
        \end{pmatrix}
        \begin{pmatrix}
        132 \\
        24 \\
        92
        \end{pmatrix}
        \]
        
        \[
        \hat{\beta} = \begin{pmatrix}
        4 \\
        -0.2 \\
        1.6
        \end{pmatrix}
        \]
        
        Así que los estimadores son:
        
        \[
        \hat{\beta_0} = 4, \quad \hat{\beta_1} = -0.2, \quad \hat{\beta_2} = 1.6
        \]
    
        \item Indique la ecuación del modelo predicho e interprete.\\

        \textbf{Respuesta:} La ecuación del modelo predicho es

        $$\hat{y} = 4 - 0.2 \cdot x_1 + 1.6 \cdot x_2$$

        Como el modelo es lineal y en niveles, esto indica que un aumento en una unidad de la variable $x_1$ hará caer a la variable $y$ en 0,2 unidades; mientras que si la variable $x_2$ varía en una unidad, la variable dependiente se incrementará en 1,6 unidades.\\

        \item ¿Qué debería suceder con $R^2$ si agregamos más variables independientes a la medición?

        \textbf{Respuesta:} El \( R^2 \) se puede escribir matemáticamente como:

    \[
    R^2 = 1 - \frac{\sum_{i=1}^{N} \hat{U_i}^2}{\sum_{i=1}^{N} (Y_i - \bar{Y})^2}
    \]

    Por lo tanto, el único valor que cambia cuando se agregan variables es \( \sum_{i=1}^{N} \hat{U_i}^2 \), que representa la función valor de la función objetivo, o sea, esta última evaluada en los parámetros que la minimizan. A su vez, por propiedad general de la optimización, la función objetivo no puede empeorar si relajamos una restricción. Por ello, la suma de los residuos al cuadrado (la función objetivo de MCO) no puede empeorar si agregamos un regresor, ya que agregar un regresor es equivalente a relajar la restricción, en el sentido de que ahora el parámetro que acompaña a ese regresor no se impone igual a cero. Finalmente, al observar la expresión inicial, una suma de residuos al cuadrado más pequeña implica un mayor \( R^2 \).



        \item ¿Qué valor alcanza el $R^2$? \footnote{Hint: Recuerde que $SSE = \sum_{i=1}^N (\hat{Y_i} - \bar{Y})^2$ que de manera matricial queda como $SSE=\hat{Y}^T Y - N \bar{Y^2}$} \\

        \textbf{Respuesta:} 
        %regresion lineal multiple

    \end{enumerate}

\end{comment}

      \item (20 puntos) Una consultora quiere medir el efecto de un programa de capacitación laboral sobre el salario mensual de los trabajadores seis meses después. Se observa la siguiente muestra de 6 personas:

    \[
    \begin{array}{|c|c|c|c|}
    \hline
    \text{Persona} & \text{Salario } (Y) & \text{Tratamiento } (T) & \text{Motivación } (M)\\
    \hline
    1 & 560 & 1 & 10\\
    2 & 520 & 1 & 8\\
    3 & 500 & 1 & 7\\
    4 & 420 & 0 & 6\\
    5 & 400 & 0 & 4\\
    6 & 380 & 0 & 5 \\
    \hline
    \end{array}
    \]

    donde:
    \begin{itemize}
        \item \(Y\): salario mensual, medido en miles de pesos.
        \item \(T=1\): el trabajador participó en la capacitación.
        \item \(T=0\): el trabajador no participó en la capacitación.
        \item \(M\): índice de motivación laboral.
    \end{itemize}

    Considere primero el siguiente modelo:
    \[
    Y_i=\beta_0+\beta_1 T_i+u_i
    \]

    Luego, considere el modelo ampliado:
    \[
    Y_i=\beta_0+\beta_1 T_i+\beta_2 M_i+u_i
    \]

    Responda:

    \begin{enumerate}
        \item (5 puntos) Calcule el estimador \(\hat{\beta}_1\) en la regresión simple de \(Y\) sobre \(T\) e interprete dentro del contexto.
        \item (5 puntos) Discuta si este estimador puede interpretarse como el efecto causal de la capacitación sobre el salario y la posibilidad de distintos sesgos en este modelo.
        \item (5 puntos) Si se incorpora la variable motivación \(M\) en la regresión, discuta si esperaría que el coeficiente estimado de \(T\) aumente, disminuya o se mantenga igual. Justifique.
        \item (5 puntos) Si al incluir la variable motivación \(M\) el modelo queda dado por:
        $$\hat{Y}_i = 345 + 90T_i + 11M_i$$
        ¿De que valor era el sesgo en este caso?
    \end{enumerate}

%     \textbf{Solución}

%     \begin{enumerate}
%         \item Como \(T\) es una variable binaria, en una regresión simple con intercepto el coeficiente \(\hat{\beta}_1\) se puede obtener como la diferencia entre el promedio de \(Y\) en el grupo tratado y el promedio de \(Y\) en el grupo no tratado.

%         Promedio del salario en el grupo tratado:
%         \[
%         \bar{Y}_1=\frac{560+520+500}{3}=\frac{1580}{3}=526.67
%         \]

%         Promedio del salario en el grupo no tratado:
%         \[
%         \bar{Y}_0=\frac{420+400+380}{3}=\frac{1200}{3}=400
%         \]

%         Por lo tanto,
%         \[
%         \hat{\beta}_1=\bar{Y}_1-\bar{Y}_0=526.67-400=126.67
%         \]

%         Entonces,
%         \[
%         \boxed{\hat{\beta}_1=126.67}
%         \]

%       La interpretación económica es que los trabajadores que participaron en la capacitación presentan, en promedio, un salario mensual mayor en \(126.67\) miles de pesos respecto de quienes no participaron.

% \item Este estimador no necesariamente puede interpretarse como el efecto causal de la capacitación. La diferencia observada entre tratados y no tratados podría deberse no solo al tratamiento, sino también a diferencias previas entre ambos grupos. \\ 
%   Puede existir sesgo de selección porque los trabajadores que participaron en la capacitación tienen niveles de motivación más altos que los no tratados. Verifiquemos esto con los promedios de motivación:

%         \[
%         \bar{M}_1=\frac{9+8+7}{3}=8
%         \qquad
%         \bar{M}_0=\frac{4+3+2}{3}=3
%         \]

%         Por lo tanto, los tratados tienen en promedio una motivación mayor. Si los trabajadores más motivados tienen más probabilidad de entrar al programa, entonces la asignación al tratamiento no es aleatoria. \\ 
%         También puede existir sesgo por variable omitida porque la motivación afecta al salario y, al mismo tiempo, está correlacionada con la variable tratamiento.

%         \item  Si la regresión omitiera \(M\), entonces esta variable quedaría contenida en el término de error \(u_i\). En ese caso, \(T_i\) estaría correlacionada con el error y se violaría el supuesto:
%         \[
%         E[u_i \mid T_i]=0
%         \]

%         Así, el coeficiente de tratamiento en la regresión simple recogería tanto el efecto de la capacitación como parte del efecto de la motivación.

%         \item Si se incorpora la variable motivación \(M\) en la regresión
%         \[
%         Y_i=\beta_0+\beta_1 T_i+\beta_2 M_i+u_i
%         \]
%         se esperaría que el coeficiente estimado de \(T\) disminuya.

%         La razón es que en la regresión simple el efecto estimado de \(T\) estaba sobrestimado, ya que capturaba:
%         \begin{itemize}
%             \item el efecto propio de la capacitación;
%             \item parte del efecto positivo de la motivación sobre el salario.
%         \end{itemize}

%         Como la motivación está positivamente correlacionada con el tratamiento y además afecta positivamente al salario, la omisión de \(M\) genera un sesgo positivo en la estimación de \(\beta_1\). Al controlar por motivación, ese sesgo debería reducirse, y por lo tanto el coeficiente de \(T\) debería bajar.


%         \item Como pasamos de un $\tilde{\beta_1}=126,67$ a $\hat{\beta_1}=90$ el valor del sesgo esta dado por:
%         $$Sesgo = \tilde{\beta_1} - \hat{\beta_1} = 126.67 - 90 = 36.67$$
%     \end{enumerate}

\item (20 puntos) Sea un modelo de regresión lineal simple en el cual, en la práctica, se observan $
Y^* = Y + u, \quad \text{donde } u \sim \mathcal{N}(0,\,1),
$ y $
X^* = X + c$, \quad \text{donde } $c \neq 0$ \text{ es una constante}. Si el modelo verdadero (sin errores de medición) es $
Y = \beta_0 + \beta_1 \, X + \varepsilon$, muestre que, al estimar el modelo usando \((Y^*,\,X^*)\) en lugar de \((Y,\,X)\), se obtiene $\hat{\beta}_1^* \;=\; \hat{\beta}_1,
\quad
\hat{\beta}_0^* \;=\; \hat{\beta}_0 \;-\;\hat{\beta}_1 \, c.
$

% \textbf{solución}
% en términos esperados ,el estimador $\hat{\beta^*_1}$ esta dado por:
% $$E(\hat{\beta^*_1})  =  \frac{Cov(X^*,Y^*)}{Var(X^*)} = \frac{Cov(X+c,Y+u)}{Var(X+c)}=\frac{Cov(X+c,\beta_0 + \beta_1X+u)}{Var(X+c)}$$
% $$\frac{Cov(X,\beta_1X)+Cov(X,u)}{Var(X)}$$
% Por supuestos de regresión lineal $Cov(X,u)=0$, por lo tanto
% $$\frac{Cov(X,\beta_1X)}{Var(X)}=\frac{\beta_1Cov(X,X)}{Var(X)}=\beta_1\frac{Var(X)}{Var(X)}=\beta_1$$
% y $E(\hat{\beta^*_0})$, esta dado por:
% $$E(\hat{\beta^*_0})=E(Y^*)-\hat{\beta^*_1}E(X^*)=E(Y+u)-\beta_1E(X+c)$$
% $$E(\hat{\beta^*_0}) = E(Y)+E(u)-\beta_1E(X)-\beta_1c$$
% Se sabe que $E(u)=0$, entonces:
% $$E(\hat{\beta^*_0}) = E(Y)-\beta_1E(X)-\beta_1c=\beta_0-\beta_1c$$

% \textbf{solución alternativa para $\beta_1$}

% $$\hat{\beta_1^*} = \frac{\sum (X_i^*-\Bar{X^*})(Y_i^*-\Bar{Y^*})}{\sum (X_i^*-\Bar{X^*})^2}$$
% $$\hat{\beta_1^*} = \frac{\sum (X_i+c-(\Bar{X}+c))(Y_i+u-(\Bar{Y}+\Bar{u}))}{\sum (X_i+c-(\Bar{X}+c))^2}$$
% $$\hat{\beta_1^*} = \frac{\sum (X_i-\Bar{X})(Y_i+u-(\Bar{Y}+\Bar{u}))}{\sum (X_i-\Bar{X})^2}$$
% $$\hat{\beta_1^*} = \frac{\sum (X_i-\Bar{X})(Y_i-\Bar{Y})}{\sum (X_i-\Bar{X})^2} + \frac{\sum (X_i-\Bar{X})(u-\Bar{u})}{\sum (X_i-\Bar{X})^2}$$
% $$\hat{\beta_1^*} = \hat{\beta_1} + \frac{\sum (X_i-\Bar{X})(u-\Bar{u})}{\sum (X_i-\Bar{X})^2}$$

% el termino $\frac{\sum (X_i-\Bar{X})(u-\Bar{u})}{\sum (X_i-\Bar{X})^2}$ es igual a 0 en terminos esperados, por lo que:
%  $$E(\hat{\beta_1^*}) = E(\hat{\beta_1})$$


\item (25 puntos) Se presenta un ejercicio que estudia la relación entre el coeficiente intelectual (CI) y el puntaje en una prueba estandarizada. Se cuenta con 7 observaciones de individuos, para los cuales se registraron:

\begin{center}
\begin{tabular}{|c|c|c|}
\hline
\textbf{Individuo} & \textbf{CI} ($X$) & \textbf{Puntaje} ($Y$)\\
\hline
1 & 95  & 40 \\
2 & 100 & 42 \\
3 & 105 & 45 \\
4 & 103 & 44 \\
5 & 109 & 46 \\
6 & 115 & 45 \\
7 & 120 & 50 \\
\hline
\end{tabular}
\end{center}

Se plantea un modelo donde el Puntaje esta dado por:
$$Puntaje = \beta_0 + \beta_1 CI + u$$

\begin{enumerate}
    \item (5 puntos) Explique por qué en este modelo es importante suponer que $E(u|X)=0$. ¿Qué pasaría si este supuesto no se cumple?
    \\
    % \textbf{solución:}
    % Si este supuesto no se cumple existiria una correlación de X con u, que indicaría que hay una variable que se esta omitiendo, por lo que es estimador estaria sesgado.

    \item (15 puntos) Calcule los estimadores del modelo, $\hat{\beta_0}$ y $\hat{\beta_1}$, que describen el comportamiento del Puntaje en función del CI.

\\

% \textbf{solución}

% \\

% Calculamos $\Bar{X}$ y $\Bar{Y}$
% $$\Bar{X} = (95+100+105+103+109+115+120)/7=106.71$$
% $$\Bar{Y}=(40+42+45+44+46+45+50)/7=44.57$$
% \begin{center}
% \begin{tabular}{|c|c|c|c|c|c|c|}
% \hline
% \textbf{Individuo} & \textbf{CI} ($X$) & \textbf{Puntaje} ($Y$) & $X_i-\Bar{X}$ & $(Y_i - \Bar{Y})$ & $(X_i-\Bar{X})(Y_i-\Bar{Y})$ & $(X_i-\Bar{X})^2$\\
% \hline
% 1 & 95  & 40 & -11.71& -4.57& 53.51& 137.12\\
% 2 & 100 & 42 & -6.71 & -2.57& 17.24& 45.02 \\
% 3 & 105 & 45 & -1.71 & 0.43 & -0.74& 2.92  \\
% 4 & 103 & 44 & -3.71 & -0.57& 2.11 & 13.76 \\
% 5 & 109 & 46 & 2.29  & 1.43 & 3.27 & 5.24  \\
% 6 & 115 & 45 & 8.29  & 0.43 & 3.56 & 68.72 \\
% 7 & 120 & 50 & 13.29 & 5.43 & 72.16& 176.62 \\
% \hline
% \end{tabular}
% \end{center}
% por tanto $\beta_1$:
% $$\hat{\beta_1} = \frac{151.11}{449.4}=0.34$$
% y $\hat{\beta_0}$ sería:
% $$\hat{\beta_0}= 44.57-0.34\cdot 106.71 = 8.29$$

    \item (5 puntos) ¿Qué representan los valores de $\hat{\beta_0}$ y $\hat{\beta_1}$?\\
    % \textbf{solución}
    % \\
    % El valor de $\hat{\beta_0}$ representa el intercepto de la RLS, mientras que el valor de $\hat{\beta_1}$ indica cuanto cambia Y, cuando X cambia en 1 unidad.
\end{enumerate}

\end{enumerate}


\section{Formulario}
\begin{itemize}

\item El intervalo de confianza de $100(1-\alpha)\%$ para la media poblacional $\mu$ cuando $\sigma$ es \emph{\textbf{conocido}}:
$\color{red}{\bar{x} \pm z_{\frac{\alpha}{2}}\frac{\sigma}{\sqrt{n}}}$

\item El intervalo de confianza de $100(1-\alpha)\%$ para la diferencia de medias
$\mu_x-\mu_y$ es:
${\color{red}(\bar{X}-\bar{Y}) \pm z_{\frac{\alpha}{2}}\sqrt{\frac{\sigma_x^2}{n_x}+\frac{\sigma_y^2}{n_y}}}$

\item  Estimadores MCO.
{\color{red}
$$\hat{\beta}_0 = \bar{y} - \hat{\beta}_1 \bar{x} \quad \quad   \hat{\beta}_1 = \frac{\sum_{i=1}^{n}(x_i - \bar{x})(y_i-\bar{y})}{\sum_{i=1}^{n}(x_i - \bar{x})^2}$$}
    
\item 
	{\color{red}	\[ SST = \sum_{i=1}^{n} (y_i - \bar{y})^2 \quad \quad  SSE = \sum_{i=1}^{n} (\hat{y}_i - \bar{y})^2 \quad \quad SSR = \sum_{i=1}^{n} \hat{u}_i^2  \]}
		
		
    \item Vector de estimadores:    
{\color{red}    $\hat{\beta}=(X^TX)^{-1}X^Ty$}

\item {\color{red}\[
A = \begin{bmatrix}
a & b \\
c & d
\end{bmatrix}
\qquad
A^{-1} = \frac{1}{det(A)} \begin{bmatrix}
d & -b \\
-c & a
\end{bmatrix}
\]}
\end{itemize}

\begin{figure}[H]
    \centering
    \includegraphics[scale=0.65]{Figuras/normal.png}
\end{figure}
\end{document}